\section{Architecture and domain design}

\subsection{Layered structure}
Questory uses feature-first modules with a shared domain and contract layer. The
presentation layer depends on domain models and interfaces. Application classes
coordinate state transitions and calculations. Data adapters implement SQLite,
bundled JSON, filesystem, GPS, camera, export, Android sharing, and optional
Supabase behavior. Dependency assembly occurs once at application startup.

\begin{figure}[htbp]
\centering
\begin{tikzpicture}[
  layer/.style={draw=QuestInk,rounded corners=2mm,minimum width=0.88\textwidth,
    minimum height=13mm,align=center,font=\small},
  arr/.style={-{Latex[length=2mm]},thick,color=QuestTeal},node distance=5mm]
  \node[layer,fill=QuestCoral!18] (ui) {Presentation: Explore, Tracker, Camera, Summary, Journey, Story Studio};
  \node[layer,fill=QuestYellow!22,below=of ui] (app) {Application: run lifecycle, quest evaluation, achievements, story editing and export};
  \node[layer,fill=QuestTeal!14,below=of app] (domain) {Domain and contracts: immutable models plus repository/device interfaces};
  \node[layer,fill=QuestPaper,below=of domain] (data) {Adapters: bundled JSON, SQLite, GPS, camera/files, PNG renderer, Android share, optional Supabase};
  \draw[arr] (ui) -- (app);
  \draw[arr] (app) -- (domain);
  \draw[arr] (data) -- (domain);
\end{tikzpicture}
\caption{Layered dependency direction. Data adapters implement contracts; domain code does not depend on plugins.}
\end{figure}



This separation prevents the interface from opening databases or manipulating
files directly. It also lets widget and application tests replace device services
with deterministic fakes. Optional networking cannot change the outcome of a
local run because it is represented by a nullable contract implementation rather
than embedded in the core model.

\subsection{Shared domain language}
\begin{table}[htbp]
\centering
\small
\begin{tabularx}{\textwidth}{@{}p{31mm}X@{}}
\toprule
\textbf{Model} & \textbf{Role} \\
\midrule
\code{GeoPoint} & Latitude, longitude, UTC timestamp, and optional accuracy or altitude \\
\code{DestinationPack} & City metadata, schema and pack version, routes, quests, and discovery points \\
\code{RoutePlan} & Expected polyline, distance, duration, landmarks, safety notes, and quest IDs \\
\code{Quest} & Prompt, point, radius, caption rule, evidence expectation, and fallback \\
\code{QuestEvidence} & Retained photo path, point, timestamp, caption, and quest identity \\
\code{RunSession} & Recoverable immutable snapshot of the active lifecycle and accumulated state \\
\code{RunSummary} & Final track, active duration, distance, pace inputs, evidence, and landmarks \\
\code{StoryDocument} & Serializable 1080 by 1920 canvas, ordered elements, transforms, and source run \\
\code{Achievement} & Stable criterion, progress, target, and optional UTC unlock timestamp \\
\bottomrule
\end{tabularx}
\caption{Principal shared domain models.}
\end{table}



Stable string identifiers connect records without exposing database row numbers.
Timestamps are stored in UTC and displayed in local time. Distances stay in
meters until presentation. Models serialize to plain maps and lists so the
database and tests can round-trip them without widget dependencies.

\subsection{Dependency assembly}
The production entry point first initializes Flutter, then opens the local
database and constructs repositories and platform services. Only after this
future succeeds does it render the main navigation. Startup loading and failure
screens prevent a partially initialized UI. Supabase URL and anonymous key are
read from compile-time definitions; when either is absent, the live-share service
is simply unavailable and the complete offline application continues normally.

\subsection{Offline-first boundary}
The local path includes bundled destination content, SQLite records, application-
controlled photos, metrics, achievements, Story Studio documents, export, and
the Android share sheet. The remote path contains only anonymous route sharing.
This boundary is architectural rather than cosmetic: a network failure cannot
roll back or invalidate local state.

